% ---- pg-017 ----
\[
\bm{X} =
\begin{bmatrix}
42 & 4 \\
52 & 5 \\
48 & 4 \\
58 & 3
\end{bmatrix}
\quad\blacksquare
\]

以阵列格式考虑数据,简化了对问题的说明,并使数字计算以一种有序且有效的方式进行.从以下两点得到的是双重益处:(1)以阵列运算描述数字计算;(2)以计算机为实现计算的工具,现在在计算机上用多种语言及统计程序包来进行阵列运算.在第 2 章我们考虑数字阵列的变换.目前,我们只考虑它们作为展示数据方法的价值.

\mvsub{描述统计量}

一个大的数据集是很庞大的,而正是它的庞大严重地干扰了任何从中提取适当信息的企图.包含在数据中的许多信息,可以通过计算某些通称为描述统计量的概括数字进行估计.例如,算术平均值,或样本均值,是一种描述统计量,它提供了一种定位测量——即一个数据集的“中心值”.而从均值到所有数的距离的平方的平均值提供出用这些数的分布程度或变差.

我们将尽量依靠描述统计量来测量定位、变差以及线性结合.这些量的正式定义如下:

设 $x_{11},x_{21},\cdots,x_{n1}$ 是第一个变量的 $n$ 个测量值,则这些测量值的算术平均值是
\[
\bar{x}_1 = \frac{1}{n}\sum_{j=1}^{n} x_{j1}
\]

如果这 $n$ 个测量值代表被观测的全部测量值集合的一个子集,则 $\bar{x}_1$ 也称为第一个变量的样本均值.我们采用这个术语是由于本书的大部分内容是用于研究为了分析来自较大聚集的测量值样本而设计的计算方法.

样本均值可用 $p$ 个变量中的每一个变量的 $n$ 个测量值计算出来,因此,一般的,将有 $p$ 个样本均值:
\phantomsection\label{eq:1-1}
\begin{equation}
\bar{x}_k = \frac{1}{n}\sum_{j=1}^{n} x_{jk} \qquad k=1,2,\cdots,p \tag{1-1}
\end{equation}

分布的程度由样本方差给出,对第一个变量的 $n$ 个观测值定义为
\[
s_1^2 = \frac{1}{n}\sum_{j=1}^{n}(x_{j1}-\bar{x}_1)^2
\]

其中 $\bar{x}_1$ 是 $x_{j1}$ 的样本均值.一般的,对于 $p$ 个变量,我们有
\phantomsection\label{eq:1-2}
\begin{equation}
s_k^2 = \frac{1}{n}\sum_{j=1}^{n}(x_{jk}-\bar{x}_k)^2 \qquad k=1,2,\cdots,p \tag{1-2}
\end{equation}

有两点说明.首先许多作者定义样本方差时,用 $n-1$ 而不是 $n$ 作为除数.后面我们将看到这样做是有理论上的理由的,并且当测量值的数目 $n$ 较小时,这样做尤其适当.样本方差的两种形式总能用适当的式子区分开来.

其次,尽管记号 $s^2$ 是样本方差传统的表示方法,我们仍将考虑一个样本方差位于主对角线的数量阵列.在这种情况下,为了表明方差在阵列中的位置,使用双下标是比较方便的.因此,引入记号 $s_{ii}$ 来表示由第 $i$ 个变量的测量值计算出的方差,并有式子
\phantomsection\label{eq:1-3}
\begin{equation}
s_k^2 = s_{kk} = \frac{1}{n}\sum_{j=1}^{n}(x_{jk}-\bar{x}_k)^2 \qquad k=1,2,\cdots,p \tag{1-3}
\end{equation}

样本方差的平方根 $\sqrt{s_{ii}}$,称为样本标准差.这个变差的度量单位和观测值的单位相同.

考虑变量 1 和 2 的 $n$ 对测量值:
